\chapter{资本、政策、标准与区域格局}
\label{chap:policy-regional-landscape}

\section{先区分四类文本}

资本提供研发、库存和工厂爬坡的时间，却不证明技术成熟。政策目标表达资源方向；法律法规产生适用范围内的义务；标准提供设计、测试或管理基线，是否强制取决于引用和合同。以下比较截止 2026 年 7 月 14 日，后续适用日期不写成已经生效的义务。

\begin{figure}[H]
\centering
\small
\begin{tabular}{c@{\ $\rightarrow$\ }c@{\ $\rightarrow$\ }c@{\ $\rightarrow$\ }c}
\fbox{政策/资金方向} & \fbox{研发与场景计划} & \fbox{产品验证/标准} & \fbox{法律义务与市场监督}
\end{tabular}
\caption{政策到准入不是单向自动转化；每一步需要责任主体和证据。}
\label{fig:policy-to-market-access}
\end{figure}

\section{中国：供应链、公共平台与场景牵引}

工信部 2023 年《人形机器人创新发展指导意见》覆盖“大脑、小脑、肢体”、整机、关键部件、软件、示范场景、标准和治理，并提出 2025/2027 阶段目标\cite{miit2023humanoid}。它是产业指导文件，不是对单一产品的认证。

2026 年工信部与国务院国资委的专项行动把工业、服务和特种领域的真实任务单元列为实景实训载体，强调常态化部署和经济可行性\cite{miit2026fieldtraining}。这一变化把资源从“做出样机”进一步拉向场景、数据和迭代，但企业仍需分别满足机械、电气、网络、数据和行业监管要求。

\section{美国：创新基础设施、行业规则与自愿框架}

美国 2025 AI Action Plan 将政策分为加速创新、建设 AI 基础设施和国际输出三支柱\cite{whitehouse2025aiaction}，并非机器人产品的统一准入法。NIST AI RMF 1.0 是自愿、非行业特定的风险管理框架\cite{nist2023airmf}；其 Govern--Map--Measure--Manage 可用于组织证据，但不能替代机械安全认证、职业安全、产品责任或医疗器械等行业义务。

因此，美国路径更像“通用 AI/算力政策 + 行业监管 + 标准与责任诉讼”的组合。企业需要按使用场景确定主管机构和合同责任，不能把遵循自愿框架写成获得政府批准。

\section{欧盟：AI 与机械产品规则叠加}

欧盟 AI Act 已于 2024 年生效，但采用分阶段适用：部分条款已适用，主要规则计划自 2026 年 8 月 2 日起适用，部分嵌入受监管产品的义务更晚\cite{eu2024aiact}。截至本章日期，“即将适用”与“已经适用”必须区分。机器人是否属于高风险 AI 需根据用途、角色和附录范围判断，不能因使用神经网络自动归类。

Machinery Regulation (EU) 2023/1230 主要自 2027 年 1 月 20 日起适用\cite{eu2023machinery}。具身系统进入欧盟市场时，AI 风险、机械危害、网络安全、技术文档、合格评定和上市后监测可能叠加；模型更新若改变安全功能或预期用途，还会触发配置与责任问题。

\section{日本与韩国：实施环境、联盟和标准化}

日本 METI 2026 AI Robotics Strategy 把机器人友好环境、用户导入支持、研发/社会实施枢纽和分行业路线图放在一起\cite{meti2026robotics}。它反映日本从强机器人制造基础向 AI、数据和实际导入补链，评价成效仍需看部署、劳动力替代质量和维护网络。

韩国 MOTIE 2024 融合机器人国际标准战略计划推进国内标准和国际提案，覆盖康复与家庭服务等方向\cite{motie2024standards}；2025 K-Humanoid Alliance 则把整机、零部件、研究机构和基础模型开发组织起来\cite{motie2025humanoid}。联盟和谅解备忘录是协作基础，不等于采购、认证或量产。

\begin{table}[H]
\begingroup\hbadness=10000
\centering\small
\caption{区域路径的同口径比较（截至 2026-07-14）}
\label{tab:regional-policy-matrix}
\begin{tabularx}{\textwidth}{p{0.1\textwidth}YYYY}
\toprule
区域 & 主要抓手 & 对研发/交付的直接影响 & 准入结构 & 关键不确定性 \\
\midrule
中国 & 指导意见、公共平台、实景实训 & 场景、供应链、数据与中试资源 & 通用安全叠加行业监管 & 目标到可审计交付的转化 \\
美国 & 创新/算力政策、行业监管、自愿框架 & 资本与基础设施强，按行业合规 & 多机构、标准、合同与责任体系 & 机器人专门规则与州差异 \\
欧盟 & AI Act + Machinery Regulation & 文档、风险、更新与上市后责任前置 & 横向 AI 与产品法规叠加 & 分期适用和分类解释 \\
日本 & AI 机器人战略、友好环境、行业路线 & 强调用户导入与系统实施 & 既有产品/行业规则与标准 & 数据/软件补链和扩散速度 \\
韩国 & 产业联盟、基础模型与国际标准 & 联合研发、部件和标准提案 & 产品安全与标准体系 & 联盟成果到客户部署 \\
\bottomrule
\end{tabularx}
\endgroup
\end{table}

\section{资本与政策如何进入技术判断}

资本和政策只在改变约束时进入路线判断：是否提供长期数据采集、昂贵验证、模具与库存爬坡、公共测试场、采购需求或合规基础设施。融资额、估值、补贴目标和联盟成员数不进入模型能力评分。可验证的中间结果是测试设施建成、标准发布、采购完成、产品通过合格评定、部署给出暴露量和客户续约。

\section{知识小结与综述判断}

五个区域都在连接 AI 与机器人，但政策工具不同：中国偏供应链和场景牵引，美国偏创新基础设施与行业治理，欧盟以前置规则叠加产品法规，日本偏实施环境，韩国偏联盟和标准。综述判断是：区域竞争最终落在数据、部件、制造、验证、部署和责任体系的闭合速度；政策文本是约束和资源的证据，不是技术胜负结论。

